% SOURCE_STATEMENT: "Can Langmann–Szabo duality be used to construct renormalizable noncommutative quantum field theories?"
% Preparation and final self-audit: CONTRIBUTING.md prompts C and D; source/audit artifacts retained in the local physics_statements_2026_10_11 batch.
% !TeX program = lualatex
\documentclass[11pt]{article}
\usepackage[margin=25mm]{geometry}
\usepackage{fontspec}
\setmainfont{Latin Modern Roman}
\usepackage{amsmath,amssymb,mathtools}
\usepackage{unicode-math}
\setmathfont{Latin Modern Math}
\usepackage{xurl,hyperref}
\hypersetup{colorlinks=true,urlcolor=blue}
\setcounter{secnumdepth}{0}
\setlength{\parindent}{0pt}
\setlength{\parskip}{5pt}
\setlength{\emergencystretch}{3em}
\newcommand{\sref}[2]{\hyperlink{source:#1:#2}{[S#2]}}
\newcommand{\cataloguescope}[1]{\paragraph{Recorded target.}#1}
\newcommand{\R}{\mathbb R}
\newcommand{\C}{\mathbb C}
\newcommand{\N}{\mathbb N}
\newcommand{\Z}{\mathbb Z}
\newcommand{\Q}{\mathbb Q}
\newcommand{\F}{\mathbb F}
\newcommand{\eps}{\varepsilon}
\begin{document}
% UnsolvedMath ID 30000342; source catalogue code OWR-1109-002
\setcounter{section}{30000341}
\section{Renormalization through Langmann–Szabo Duality}\label{30000342}
{\small UnsolvedMath ID 30000342 · Mathematical Physics}
\subsection{Definitions and mathematical statement}

\paragraph{Shared notation.}$\mathbb N=\{1,2,\ldots\}$; the standard number fields have their usual meanings. All additional conventions are stated locally.


\paragraph{Definitions and conventions.} The Moyal deformation replaces ordinary multiplication of fields by a noncommutative star product with coordinate commutators $[x^\mu,x^\nu]=i\theta^{\mu\nu}$. Langmann--Szabo duality relates position and momentum descriptions of certain such scalar models, exchanging the kinetic and harmonic-oscillator terms. Perturbative renormalizability means that ultraviolet divergences at all perturbative orders can be absorbed into a fixed finite set of parameters and field normalizations.

\paragraph{Source scope and qualifications.} The original report already describes the successful harmonic scalar construction. Its broad construction question is historical and model dependent; it must not be converted into a claim that no example is known, or into an unsupported theorem for all noncommutative field theories.

\paragraph{Catalogue mathematical target.}
\paragraph{Catalogue wording retained for traceability.}
Can Langmann–Szabo duality be used to construct renormalizable noncommutative quantum field theories?

\paragraph{Saved literature context.} The supplied catalogue's assessment dated 2026-08-21 records: Langmann--Szabo duality led to renormalizable harmonic noncommutative scalar theories, while the general construction program and gauge-theory extension remain model dependent. The assessment is retained as dated provenance; it is not a new comprehensive literature or open-status certification.
\subsection{Short English statement}
Use position--momentum duality to construct and understand renormalizable noncommutative quantum field models.
\subsection{Sources}
\begin{itemize}
\item[\hypertarget{source:30000342:1}{S1}] ULAM.AI and the UnsolvedMath Contributors, UnsolvedMath, record 30000342 (OWR-1109-002). Catalogue curation under CC BY 4.0; definitions and English exposition added locally.
\url{https://huggingface.co/datasets/ulamai/UnsolvedMath}.
\item[\hypertarget{source:30000342:2}{S2}] Noncommutative Geometry and Quantum Field Theory (2005), Oberwolfach Reports, DOI 10.4171/owr/2005/48, pp. 2705--2760; official report source and mathematical context.
\url{https://doi.org/10.4171/owr/2005/48}.
\item[\hypertarget{source:30000342:3}{S3}] Official EMS Press PDF of the same Oberwolfach report; PDF pages 17--19.
\url{https://ems.press/content/serial-article-files/46022}.
\item[\hypertarget{source:30000342:4}{S4}] Langmann--Szabo duality; reference in the saved dated catalogue literature review.
\url{https://arxiv.org/abs/hep-th/0202039}.
\item[\hypertarget{source:30000342:5}{S5}] Renormalizable UV--IR duality models; reference in the saved dated catalogue literature review.
\url{https://arxiv.org/abs/1001.3776}.
\end{itemize}
\label{end:30000342}
\end{document}
